% !TEX root = ../../../main.tex
% Sources: LEC01, IMG_9114--IMG_9117.
\section{Power sets and the size of the real line}
\label{sec:analysis-i:power-sets}

The \emph{power set} $\mathcal P(X)$ is the set of all subsets of $X$.
For a finite set with $n$ elements, a subset is specified by deciding,
independently for each element, whether to include it. Thus
\[
  |\mathcal P(X)|=2^n=\sum_{k=0}^{n}\binom{n}{k}.
\]
The smallest case distinguishes the empty set from the set containing it.
% !TEX root = ../../hw01.tex
% Homework 01, Exercise 2. Shared problem statement.
\begin{exercise}
  \label{ex:hw01:power-set}
  \solutionlink{hw01:power-set}
  List all the elements in the set $\mathcal{P}(\mathcal{P}(\varnothing))$
  and verify the cardinality formula.
\end{exercise}


Cantor's theorem extends the strict increase in size to every set.

\begin{samepage}
\begin{theorem}[Cantor]
  \label{thm:analysis-i:cantor}
  There is no surjection from $X$ onto $\mathcal P(X)$.
\end{theorem}
\begin{proof}
  Let $\varphi:X\to\mathcal P(X)$ be any map, and form
  \[
    D=\{x\in X:x\notin\varphi(x)\}.
  \]
  If $D=\varphi(d)$ for some $d\in X$, then the definition gives
  $d\in D$ if and only if $d\notin D$. This is impossible, so $D$ is
  missing from the image of $\varphi$.
\end{proof}
\end{samepage}
The map $x\mapsto\{x\}$ is injective, so the theorem is expressed as
$|X|<|\mathcal P(X)|$. In particular, $\mathcal P(\N)$ is uncountable.

\subsection{Digits as subsets}
Decimal notation is a sum of place values. For example,
\[
  \pi=3+\frac{1}{10}+\frac{4}{10^2}+\frac{1}{10^3}+\cdots.
\]
More generally, a base-$b$ expansion of a number in $[0,1]$ has the form
\[
  x=\sum_{k=1}^{\infty}\frac{a_k}{b^k},
  \qquad a_k\in\{0,1,\ldots,b-1\},\quad b\geq2.
\]
Successively subdividing an interval into $b$ equal parts supplies these
digits; the remaining interval has length $b^{-n}$ after $n$ steps.
The nested interval property below justifies the limiting description.

In base two, a digit string records a subset of $\N$ through the positions
of its ones. There is a small ambiguity, such as
$0.1000\ldots=0.0111\ldots$. Choose the terminating expansion whenever
there are two choices, and represent $1$ by $0.1111\ldots$.
This gives an injection $[0,1]\to\mathcal P(\N)$.

For an injection in the other direction, send a subset $A\subseteq\N$ to
\[
  t_A=\sum_{k\in A}\frac{2}{3^k}\in[0,1].
\]
If $A$ and $B$ first differ at position $j$, the difference $2/3^j$ there
exceeds the maximum possible opposing tail
$\sum_{k>j}2/3^k=1/3^j$. Hence $t_A\ne t_B$.
The Cantor--Bernstein theorem now gives
\[
  |[0,1]|=|\mathcal P(\N)|=2^{\aleph_0}.
\]
Here we use the set-theoretic Cantor--Bernstein theorem in its usual form,
that injections each way imply a bijection; its proof is not needed for
the subsequent analysis.

The function $x\mapsto\tfrac12(1+x/(1+|x|))$ injects $\R$ into $(0,1)$.
To see injectivity, $x/(1+|x|)$ is strictly increasing on each side of zero
by cross multiplication, and has the sign of $x$. Together with the
inclusion $[0,1]\subseteq\R$, this proves the following cardinality
identity.\marginnote{The \emph{continuum hypothesis} is the additional
assertion that there is no cardinality strictly between $\aleph_0$ and
$\mathfrak c$. The identity $|\R|=2^{\aleph_0}$ does not assume that hypothesis.}
\[
  |\R|=2^{\aleph_0}=:\mathfrak c.
\]

\subsection{Exercises on the cardinality of intervals}
The next constructions compare intervals, the real line, and a square.
Here $\aleph$ denotes $|\R|=\mathfrak c$. The first construction uses the
familiar tangent function and its inverse.

% !TEX root = ../../hw01.tex
% Homework 01, Exercise 4. Shared problem statement.
\begin{exercise}
  \label{ex:hw01:interval-to-reals}
  \solutionlink{hw01:interval-to-reals}
  Obtain a bijective map from $(0,1)$ to $\mathbb{R}$.
\end{exercise}

% !TEX root = ../../hw01.tex
% Homework 01, Exercise 5. Shared problem statement.
\begin{exercise}
  \label{ex:hw01:open-to-closed-interval}
  \solutionlink{hw01:open-to-closed-interval}
  Obtain a bijective map from $(0,1)$ to $[0,1]$.
\end{exercise}

% !TEX root = ../../hw01.tex
% Homework 01, Exercise 8. Shared problem statement.
\begin{exercise}
  \label{ex:hw01:square-cardinality}
  \solutionlink{hw01:square-cardinality}
  Show that
  \[
    |(0,1)\times(0,1)|=|(0,1)|=\aleph.
  \]
  (\textit{Corrected hint.} For each number in $(0,1)$, choose the decimal
  expansion that is not eventually all $9$s. We can then write
  \[
    x\in(0,1),\quad x=0.x_1x_2\cdots,\quad
    x_i\in\{0,\ldots,9\},\quad i=1,2,\ldots.
  \]
  Then define a map from $(0,1)\times(0,1)$ to $(0,1)$ by
  \[
    (x,y)=(0.x_1x_2\cdots,0.y_1y_2\cdots)
    \mapsto (0.x_1y_1x_2y_2\ldots)=z\in(0,1).
  \]
  Show that this map is injective, then use injections in both directions.)
\end{exercise}


\subsection{Algebraic and transcendental numbers}
A real number is \emph{algebraic} if it is a root of a nonzero polynomial
with rational coefficients. It is \emph{transcendental} otherwise.
For instance, $\sqrt{2}$ is algebraic because it satisfies $x^2-2=0$.
Rational numbers are algebraic as well, since $p/q$ satisfies $qx-p=0$.
\index{algebraic number}\index{transcendental number}

\begin{proposition}
  The set $\mathcal A$ of real algebraic numbers is countably infinite.
  The set $\mathcal T=\R\setminus\mathcal A$ has cardinality $\mathfrak c$.
\end{proposition}
\begin{proof}
  For each $m\geq1$, polynomials of degree $m$ with rational coefficients
  are encoded by a subset of $\Q^{m+1}$, which is countable by repeated
  products of countable sets. Each nonzero degree-$m$ polynomial has at
  most $m$ real roots. Its factorization by $x-r$ at a root, followed by
  induction on the degree, proves this bound. Thus the set of roots of
  all such polynomials is countable. Taking the union over $m$ shows that
  $\mathcal A$ is countable. It is infinite because it contains $\Q$.

  If $\mathcal T$ were countable, then $\R=\mathcal A\cup\mathcal T$
  would be countable, contrary to Cantor's theorem. Choose distinct
  $t_1,t_2,\ldots\in\mathcal T$, and enumerate
  $\mathcal A=\{a_1,a_2,\ldots\}$ without repetition.
  Map $a_n$ to $t_{2n}$ and $t_n$ to $t_{2n-1}$, and leave every other
  transcendental number fixed. This is a bijection from $\R$ to
  $\mathcal T$, proving the cardinality assertion.
\end{proof}

Counting proves that transcendental numbers exist, without identifying a
particular one. The lecture also named $e$, $\pi$, and the Liouville number
\[
  \ell=\sum_{n=1}^{\infty}10^{-n!}.
\]
Their transcendence is stated here without proof. Including the $n=0$
term gives $\ell+1/10$, which is transcendental
as well. These stronger results belong to transcendental number theory~\cite{Waldschmidt2017}.
The earlier proof established only the irrationality of $e$.

\subsection{The Jacobian conjecture}
\label{sec:analysis-i:jacobian-conjecture}
Another lecture question concerns polynomial maps. A map
$F=(F_1,\ldots,F_d):\C^d\to\C^d$ is \emph{polynomial} when each
coordinate $F_i$ is a polynomial in $z_1,\ldots,z_d$. Its
\emph{Jacobian matrix} is the matrix of first derivatives,
\[
  \mathrm DF(z)=
  \left(\frac{\partial F_i}{\partial z_j}(z)\right)_{1\leq i,j\leq d}.
\]

\begin{conjecture}[Jacobian]
  If a polynomial map $F:\C^d\to\C^d$ satisfies
  $\det\mathrm DF(z)=c$ for every $z\in\C^d$, where $c\ne0$ is
  constant, then it has a polynomial inverse. That is, there is a
  polynomial map $G:\C^d\to\C^d$ such that
  \[
    G\circ F=F\circ G=\operatorname{id}_{\C^d}.
  \]
\end{conjecture}
This is the classical formulation~\cite{deGoursac2016}.

The determinant condition makes the linear approximation invertible at
every point. The inverse function theorem, a result from later differential
calculus, gives $F$ an inverse on a neighborhood of each point.
The conjecture asks for one inverse on all of $\C^d$, with polynomial
coordinates.

Figure~\ref{fig:jacobian-shear} illustrates a simple case,
$F(x,y)=(x+y^2,y)$, on the real plane inside $\C^2$. Here
\[
  \mathrm DF(x,y)=\begin{pmatrix}1&2y\\0&1\end{pmatrix},
  \qquad \det\mathrm DF(x,y)=1.
\]
The inverse $G(u,v)=(u-v^2,v)$ keeps the second coordinate and undoes
the horizontal shift. This verifies the conjecture for this map; the
general problem remains open~\cite{deGoursac2016}.
\begin{figure}[htbp]
  \centering
  % A finite real grid and its exact image under F(x,y)=(x+y^2,y).
\begin{tikzpicture}[analysis diagram,x=0.9cm,y=0.9cm]
  \begin{scope}
    \foreach \x in {-0.5,0,0.5} {
      \draw[gray,line width=0.5pt] (\x,-1) -- (\x,1);
    }
    \foreach \y in {-1,-0.5,0,0.5,1} {
      \draw[gray,line width=0.5pt] (-0.5,\y) -- (0.5,\y);
    }
    \draw[->] (-1,0) -- (1.15,0) node[right] {$x$};
    \draw[->] (0,-1.25) -- (0,1.35) node[above] {$y$};
    \fill (0,0) circle[radius=1.6pt];
    \node[analysis label,below left] at (0,0) {$0$};
    \fill (0,1) circle[radius=1.6pt];
    \node[analysis label,above left] at (0,1) {$(0,1)$};
    \node at (0,-1.7) {Original grid};
  \end{scope}

  \draw[->] (1.5,0.5) -- (3.65,0.5)
    node[midway,above] {$F$};
  \draw[->] (3.65,-0.5) -- (1.5,-0.5)
    node[midway,below] {$G=F^{-1}$};

  \begin{scope}[xshift=4.6cm]
    \foreach \x in {-0.5,0,0.5} {
      \draw[gray,line width=0.5pt,domain=-1:1,samples=41,variable=\t]
        plot ({\x+\t*\t},{\t});
    }
    \foreach \y in {-1,-0.5,0,0.5,1} {
      \draw[gray,line width=0.5pt]
        ({-0.5+\y*\y},\y) -- ({0.5+\y*\y},\y);
    }
    \draw[->] (-0.9,0) -- (1.95,0) node[right] {$u$};
    \draw[->] (0,-1.25) -- (0,1.35) node[above] {$v$};
    \fill (0,0) circle[radius=1.6pt];
    \node[analysis label,below left] at (0,0) {$0$};
    \fill (1,1) circle[radius=1.6pt];
    \node[analysis label,above] at (1,1) {$(1,1)$};
    \node at (0.5,-1.7) {Image under $F$};
  \end{scope}
\end{tikzpicture}

  \caption[A polynomial shear and its inverse.]{A real slice of the
    polynomial shear $F(x,y)=(x+y^2,y)$. Horizontal grid segments stay
    horizontal; vertical ones become parabolic arcs. The inverse
    subtracts $v^2$ from the first coordinate. This example has a
    global polynomial inverse.}
  \label{fig:jacobian-shear}
\end{figure}

\newthought{Motivation.} Earlier we studied injectivity
and surjectivity; a global inverse requires both. The Jacobian determinant
tests behavior near each point, so the conjecture connects those global
questions with local differential calculus. Requiring the inverse to be
polynomial adds an algebraic constraint. It illustrates why the precise
hypotheses of a theorem matter, just as existence by counting does not
identify a particular transcendental number. None of the arguments here
depends on its resolution.

\subsection{Exercises on transcendence}
The first exercise asks for the cardinality conclusion under CH, although
the argument above establishes it without that additional hypothesis.
The second applies the definition of transcendence to linear independence.

% !TEX root = ../../hw01.tex
% Homework 01, Exercise 6. Shared problem statement.
\begin{exercise}
  \label{ex:hw01:transcendental-cardinality}
  \solutionlink{hw01:transcendental-cardinality}
  Let $\aleph_0$ and $\aleph$ denote the cardinalities of $\mathbb{N}$ and
  $\mathbb{R}$ respectively. Recall that the Continuum Hypothesis (CH)
  states that there is no cardinality between $\aleph_0$ and $\aleph$.
  Use CH to prove that the set of transcendental numbers, $\mathbb{T}$,
  has the same cardinality $\aleph$ as that of $\mathbb{R}$.
\end{exercise}

% !TEX root = ../../hw01.tex
% Homework 01, Exercise 7. Shared problem statement.
\begin{exercise}
  \label{ex:hw01:linear-independence}
  \solutionlink{hw01:linear-independence}
  Let $\tau$ be a transcendental number whose existence is proven already.
  Show that for any $n\in\mathbb{N}$, the numbers
  \[
    \tau,\ldots,\tau^n
  \]
  are linearly independent over the field $\mathbb{Q}$.

  That is, if there are $r_1,\ldots,r_n\in\mathbb{Q}$ such that
  \[
    r_1\tau+\cdots+r_n\tau^n=0,
  \]
  then $r_1=\cdots=r_n=0$.

  (In the language of Linear Algebra, this result indicates that the set
  $\mathbb{R}$ viewed as a vector space over $\mathbb{Q}$ is infinitely
  dimensional.)
\end{exercise}

